\uuid{pj4J}
\exo7id{6073}
\titre{exo7 6073}
\auteur{queffelec}
\organisation{exo7}
\datecreate{2011-10-16}
\isIndication{false}
\isCorrection{false}
\chapitre{Continuité, uniforme continuité}
\sousChapitre{Continuité, uniforme continuité}
\cpge{mpsi-mp2i.limites-continuite.b, pcsi.limites-continuite.b, ptsi.limites-continuite.b, mpsi-mp2i.applications-lineaires.a, pcsi.applications-lineaires.a, ptsi.applications-lineaires.a, bcpst1.applications-lineaires-matrices, mp-mpi.topologie-espaces-normes.e, pc.espaces-vectoriels-normes.e, psi.espaces-vectoriels-normes.e, pt.fonctions-deux-trois-variables.b, mp-mpi.topologie-espaces-normes.f, bcpst2.applications-lineaires-matrices.a}
\module{Topologie}
\niveau{L3}
\difficulte{}

\contenu{
\texte{
Soit $E$ et $F$ deux espaces vectoriels normés et on note $B_E$ la boule unité
fermée de $E$. Soit $u$ une application de $E$ dans $F$ telle que 

\quad (i) $u(x+y)=u(x)+u(y),\ \forall x,y\in E$.

\quad (ii) $u(B_E)$ est bornée dans $F$.
}
\begin{enumerate}
    \item \question{Calculer $u(rx)$, $x\in E$, $r$ rationnel.}
    \item \question{Montrer que $u$ est continue en $0$, plus précisément : 

$$\exists M>0\ ;\ \forall x\not=0\ \ ||u(x)||\leq M||x||.$$}
    \item \question{Montrer que $u$ est continue et linéaire.}
\end{enumerate}
}
